\section{Edge Detection Essentials}
\label{app:edges}

Every voting method in Section~\ref{sec:visual-object-models} starts from edge
pixels and their gradients. The intuition is simple: an edge is where the
brightness changes abruptly, so its location is where the first derivative of
the image is large. This appendix collects what is needed to compute gradients
and to understand the reference edge detector, Canny.

\paragraph{Gradients from derivative filters.}
On a pixel grid the derivatives $G_x = \partial I/\partial x$ and
$G_y = \partial I/\partial y$ are approximated by small convolution kernels. The
Sobel kernels combine a central difference in one direction with a smoothing
in the other:
\begin{equation}
  S_x = \begin{pmatrix} -1 & 0 & 1\\ -2 & 0 & 2\\ -1 & 0 & 1 \end{pmatrix},
  \qquad
  S_y = S_x^\top = \begin{pmatrix} -1 & -2 & -1\\ 0 & 0 & 0\\ 1 & 2 & 1
  \end{pmatrix}.
\end{equation}
\noindent
The Prewitt kernels are the same with the weights $2$ replaced by $1$. From the
two responses follow the gradient magnitude
$\lVert\nabla I\rVert = \sqrt{G_x^2 + G_y^2}$ (edge strength) and the orientation
$\theta = \operatorname{atan2}(G_y,G_x)$. The gradient points from dark to
bright and is perpendicular to the edge; the edge itself runs along
$\theta \pm 90^\circ$.

\paragraph{The Canny detector.}
Thresholding the gradient magnitude alone gives thick, broken, noisy edges.
Canny fixes this in four steps:
\begin{enumerate}
  \item smooth the image with a Gaussian to suppress noise;
  \item compute gradient magnitude and orientation, e.g.\ with Sobel kernels;
  \item \emph{non-maximum suppression}: keep a pixel only if its magnitude is a
        local maximum along the gradient direction, which thins every edge to a
        one-pixel-wide line;
  \item \emph{hysteresis thresholding} with a high and a low threshold: pixels
        above the high threshold are strong edges and are kept; pixels between
        the thresholds are weak edges and are kept only if they connect to a
        strong edge; everything below the low threshold is discarded.
\end{enumerate}
\noindent
Hysteresis removes isolated noise responses while preserving long, coherent
boundaries. The detector fails where the image offers no contrast: low-contrast
boundaries, haze, and gradual transitions stay below threshold, and strong
texture produces fragmented contours. The smoothing width and the two thresholds
must be tuned to the imaging conditions.

\section{Fourier Transform and Fourier Descriptors}
\label{app:fourier}

The Fourier transform rewrites a signal as a sum of sinusoids of different
frequencies; the list of their complex weights is the \emph{spectrum}. Low
frequencies describe the coarse shape of a signal, high frequencies its fine
detail and noise. Two applications appear in object models: correlating
gradient images quickly as an alternative to R-tables, and describing a closed
contour by the spectrum of its boundary (Fourier descriptors). Both need only the
facts collected here.

\subsection{The Discrete Fourier Transform}

For a discrete signal $x(0),\dots,x(N-1)$ the discrete Fourier transform (DFT)
and its inverse are
\begin{equation}
  X(k) = \sum_{n=0}^{N-1} x(n)\, e^{-2\pi i kn/N},
  \qquad
  x(n) = \frac{1}{N}\sum_{k=0}^{N-1} X(k)\, e^{2\pi i kn/N}.
\end{equation}
\noindent
The spectrum has as many coefficients as the signal has samples, and it is
periodic, $X(k) = X(k+N)$, so the upper half of the indices can equally be read
as negative frequencies: $X(N-1) = X(-1)$, $X(N-2) = X(-2)$, and so on.
Implementations therefore often store the spectrum as
$[X(0),X(1),\dots,X(N/2),X(-N/2),\dots,X(-2),X(-1)]$. The zero-frequency term
$X(0) = \sum_n x(n)$ is $N$ times the mean of the signal.

\paragraph{Real signals.}
If $x$ is real-valued, the spectrum is \emph{Hermitian}: $X(-k) = X(k)^*$,
where $^*$ is complex conjugation. Half of the spectrum then determines the
other half. For example, if a real signal with $N=17$ samples has
$X(0),\dots,X(8) = [10,\ -5+i,\ -0.5+0.2i,\ -0.04+0.03i,\ 0,0,0,0,0]$, the
remaining entries $X(-8),\dots,X(-1)$ are
$[0,0,0,0,0,\ -0.04-0.03i,\ -0.5-0.2i,\ -5-i]$: conjugate and mirror, without
changing the sign of the real parts.

\paragraph{Standard spectra.}
Table~\ref{tab:fourier-pairs} lists the transform pairs used most often.

\begin{table}[ht]
  \centering
  \begin{tabular}{@{}p{4.2cm}p{8.8cm}@{}}
    \toprule
    Signal & Spectrum \\
    \midrule
    $\cos(2\pi f_0 x)$ & two real spikes at $\pm f_0$ \\
    $\sin(2\pi f_0 x)$ & two imaginary spikes at $\pm f_0$ with opposite signs \\
    rectangle of width $w$ & sinc function, $\propto \operatorname{sinc}(wf)$;
      a narrower rectangle gives a wider sinc \\
    sinc function & rectangle (ideal low-pass) \\
    Gaussian & Gaussian; a narrow Gaussian gives a wide one \\
    \bottomrule
  \end{tabular}
  \caption{Common Fourier pairs. Width in the signal domain and width in the
  frequency domain are inversely related.}
  \label{tab:fourier-pairs}
\end{table}

\paragraph{Convolution and correlation theorems.}
Convolution in the signal domain is multiplication in the frequency domain, and
correlation is the same with \emph{one} factor conjugated:
\begin{equation}
  f * g = \mathrm{DFT}^{-1}\bigl(\mathrm{DFT}(f)\cdot\mathrm{DFT}(g)\bigr),
  \qquad
  f \star g = \mathrm{DFT}^{-1}\bigl(\mathrm{DFT}(f)^*\cdot\mathrm{DFT}(g)\bigr).
  \label{eq:corr-theorem}
\end{equation}
\noindent
Conjugating both factors, or neither, does not give the correlation. The
practical value is speed: a correlation of a template with a whole image costs
two forward DFTs, a pointwise product, and one inverse DFT.

\paragraph{Gradient correlation instead of R-tables.}
Write the gradients of the image $I$ and of the template $O$ as complex numbers,
\begin{equation}
  I_I = \frac{\partial I}{\partial x} + i\,\frac{\partial I}{\partial y},
  \qquad
  O_I = \frac{\partial O}{\partial x} + i\,\frac{\partial O}{\partial y}.
\end{equation}
\noindent
For two complex numbers $a = a_x + i a_y$ and $b = b_x + i b_y$ the product
$a^* b = (a_x b_x + a_y b_y) + i\,(a_x b_y - a_y b_x)$ has the dot product of the
two gradient vectors as its real part. Hence
$\operatorname{Re}(I_I \star O_I)(x,y)$ sums, over the whole template placed at
$(x,y)$, the dot products between template gradients and image gradients. It is
large where many image edges have the position and orientation the template
expects, which is almost the same as Hough voting with an R-table. The
imaginary part contains cross products instead and is not the voting score.
Equation~\eqref{eq:corr-theorem} computes the correlation for all positions at
once.

\subsection{Fourier Descriptors}

A closed contour can be described as a periodic complex signal. Trace the
boundary of a segmented object at $N$ points and write each point as one
complex number,
\begin{equation}
  u(n) = x(n) + i\,y(n), \qquad n = 0,\dots,N-1 .
\end{equation}
\noindent
The Fourier descriptor is its DFT, $F(\mu) = \mathrm{DFT}(u)(\mu)$ for
$\mu = 0,\dots,N-1$. From a greyscale image the pipeline is therefore: segment the
object, trace its boundary at $N$ points, form $u$, apply the DFT, normalize,
and keep the low-frequency coefficients. Two counting facts follow directly
from the DFT. The descriptor has exactly as many coefficients as the contour has
points, so a descriptor $F(\mu)$ with $\mu\in\{0,\dots,69\}$ describes a contour
of $70$ points, whatever the value of $F(0)$. And
$F(0) = \sum_n u(n) = N\cdot(\bar x + i\,\bar y)$ is $N$ times the centroid: for
a contour of $N = 20$ points centered at $(10,20)$, $F(0) = 200 + 400i$ (with a
$1/N$-normalized DFT it would be the centroid $10 + 20i$ itself).

\paragraph{How transformations act.}
Each geometric change of the shape has a simple effect on the spectrum, which is
what makes normalization possible (Table~\ref{tab:fd-norm}). A translation by
$t = t_x + i t_y$ adds $t$ to every $u(n)$ and therefore only changes $F(0)$ (by
$Nt$). A scaling by $a$ multiplies every coefficient by $a$. A rotation by
$\varphi$ multiplies every coefficient by $e^{i\varphi}$, a pure phase change. A
different start point on the contour multiplies $F(\mu)$ by a phase
$e^{2\pi i\mu m/N}$ that grows with $\mu$.

\begin{table}[ht]
  \centering
  \begin{tabular}{@{}p{2.8cm}p{4.3cm}p{5.9cm}@{}}
    \toprule
    Invariance to & Effect on $F$ & Normalization \\
    \midrule
    Translation & changes only $F(0)$ & set $\tilde F(0) := 0$ \\
    Scale & multiplies all $F(\mu)$ by $a$ & divide by $\lvert F(k)\rvert$ for a
      fixed $k\neq 0$, e.g.\ $k=1$ \\
    Rotation & multiplies all $F(\mu)$ by $e^{i\varphi}$ & keep magnitudes
      $\lvert F(\mu)\rvert$ (or normalize the phase) \\
    Start point & phase $e^{2\pi i \mu m/N}$ & keep magnitudes
      $\lvert F(\mu)\rvert$ \\
    \bottomrule
  \end{tabular}
  \caption{Normalizing a Fourier descriptor. Each step removes one nuisance
  transformation; mirroring and noise are not removed by any of them.}
  \label{tab:fd-norm}
\end{table}

\noindent
The normalization $\tilde F(0) := 0$, $\tilde F(\mu) := F(\mu)/\lvert F(2)\rvert$
therefore makes the descriptor invariant to translation (first step) and to
scale (second step, with $k=2$). It is not invariant to rotation, because the
phases are kept. It is not invariant to mirroring, since mirroring conjugates
$u$ and exchanges positive and negative frequencies. Noise is never removed by
normalization. Plots of normalized descriptors can be recognized by the same
rules: the zero-frequency value is zero after translation normalization, the
reference coefficient has magnitude $1$ after scale normalization, and all
values are real and non-negative after taking magnitudes.

\paragraph{Smoothing a shape.}
Fine detail and noise of the contour live in the high frequencies, which are the
indices in the middle of the stored array (large $\lvert\mu\rvert$). A shape is
smoothed by a low-pass operation: set those coefficients to zero or attenuate
them, keep the low frequencies at both ends of the array, and apply the inverse
DFT. Keeping only magnitudes destroys the shape, because the phases carry the
positions of the points. Dividing the whole spectrum by a constant only rescales
the shape and does not smooth it.

\paragraph{Limitations.}
Fourier descriptors need a clean segmentation and one closed, connected contour.
They are global: a partial occlusion, or two touching objects, changes the whole
contour and therefore every coefficient. The number of contour points and the
sampling must be consistent across shapes to compare descriptors.

\section{Classical Classifiers and Part-Based Detectors}
\label{app:classifiers}

The classification stage of the processing chain in
Figure~\ref{fig:processing-chain} can be filled by many learners. Part-based
detectors build on them. This appendix covers the classifiers and two
detectors that exercises ask about beyond the implicit shape model: the Hough
forest and the deformable part model. For each, the aim is to know its working
principle and what exactly it outputs.

\paragraph{Linear classifiers and SVMs.}
A linear classifier decides by the sign of an affine function,
$g(x) = w^\top x + b$; its decision boundary is a hyperplane. For linearly
separable data many hyperplanes classify the training set correctly. A
\emph{linear support vector machine} (SVM) picks the one with the
\emph{maximum margin}: the hyperplane whose distance to the nearest training
samples of both classes is largest, running through the middle of the gap. Only
these nearest samples, the support vectors, determine it. A soft-margin variant
tolerates some misclassified samples, and kernels allow nonlinear boundaries.
In a figure with several separating lines, the SVM line is the one that keeps
equal and largest distance to both classes, not one that touches a sample.

\paragraph{Random forests.}
A decision tree routes a sample through binary node tests (for example ``is
feature $j$ larger than $\tau$?'') to a leaf that stores a class distribution.
During training, each node picks, from a set of randomly generated candidate
tests, the one that best reduces the class entropy of the samples it splits
(information gain). A random forest trains many trees, each on a random subset
of the data and with random candidate tests, and averages the leaf
distributions of all trees. The randomness makes the trees differ, so their
average has lower variance than any single tree.

\paragraph{Hough forests.}
A Hough forest is a random forest whose trees vote like the implicit shape
model. Training samples are image patches, each labeled as object or background
and, for object patches, with the offset vector to the object center. The node
tests are simple and fast: compare the values of one feature channel (intensity,
gradient, and so on) at two pixel positions inside the patch,
$f(p) < f(q) + \tau$. Each node is optimized for one of two energies, chosen at
random: the \emph{class-label uncertainty} (entropy of object versus background)
or the \emph{offset uncertainty} (variance of the offset vectors of the object
patches). Each leaf stores the fraction of object patches that reached it (a
class probability) and the list of their offset vectors. At detection time every
patch of the image descends all trees, and each leaf casts its offsets as votes
for the object center, weighted by its class probability. The codebook and its
matching step are thus replaced by a learned mapping from appearance to votes.

\paragraph{Deformable part models.}
A deformable part model (DPM) describes an object by a coarse \emph{root}
filter on HOG features, several \emph{part} filters at twice the resolution,
and a quadratic deformation cost for each part's displacement from its anchor
position. The score of a hypothesis with root position $p_0$ is
\begin{equation}
  \mathrm{score}(p_0) = F_0\cdot\phi(p_0)
  + \sum_{i=1}^{n} \max_{p_i}\Bigl[F_i\cdot\phi(p_i)
  - d_i\cdot\bigl(dx_i,\,dy_i,\,dx_i^2,\,dy_i^2\bigr)\Bigr] + b,
\end{equation}
\noindent
where $\phi(p)$ is the HOG feature window at position $p$, $F_i$ are the filter
weights, and $(dx_i,dy_i)$ is part $i$'s displacement from its anchor. The part
positions are latent variables, and the filters are trained with a latent SVM.
Localization is a sliding-window search: the score is evaluated at every
position and scale of a HOG feature pyramid, and the high-scoring windows,
after non-maxima suppression, are the detections. The SVM only outputs a score
for a given window, not a location. Objects are found by classifying windows at
many locations, not by a codebook and Hough voting.
